\documentclass[10pt]{beamer}
\usetheme{Warsaw}
\usepackage[T1]{fontenc}
\usepackage[utf8]{inputenc}
\usepackage{listings}
\usepackage{xcolor}
\usepackage{listingsutf8}

% Configuration pour le Markdown
\lstset{
  inputencoding=utf8,
  basicstyle=\ttfamily\small,
  upquote=true,
  literate=%
         {á}{{\'a}}1
         {à}{{\`{a}}}1
         {í}{{\'i}}1
         {é}{{\'e}}1
         {ý}{{\'y}}1
         {ú}{{\'u}}1
         {ó}{{\'o}}1
         {ě}{{\v{e}}}1
         {š}{{\v{s}}}1
         {č}{{\v{c}}}1
         {ř}{{\v{r}}}1
         {ž}{{\v{z}}}1
         {ď}{{\v{d}}}1
         {ť}{{\v{t}}}1
         {ň}{{\v{n}}}1                
         {ů}{{\r{u}}}1
         {Á}{{\'A}}1
         {Í}{{\'I}}1
         {É}{{\'E}}1
         {Ý}{{\'Y}}1
         {Ú}{{\'U}}1
         {Ó}{{\'O}}1
         {Ě}{{\v{E}}}1
         {Š}{{\v{S}}}1
         {Č}{{\v{C}}}1
         {Ř}{{\v{R}}}1
         {Ž}{{\v{Z}}}1
         {Ď}{{\v{D}}}1
         {Ť}{{\v{T}}}1
         {Ň}{{\v{N}}}1                
         {Ů}{{\r{U}}}1    
}

\begin{document}

\begin{frame}[fragile]
  \frametitle{Slipshow}

  \begin{itemize}
    \item Slipshow est un logiciel libre pour faire des présentations améliorées. \pause
    \item Avec slipshow, pas besoin de gérer l'alignement du texte ! \pause
    \item Une présentation slipshow prend la forme d'un fichier texte.
  \end{itemize}

  \pause

  \vspace{0.25cm}

  \textbackslash includegraphics[width=0.6\textbackslash textwidth,height=0.4\textbackslash textwidth]\{slipshow.png\}

  \vspace{0.25cm}

  \pause

  \begin{exampleblock}{Example}
    \begin{lstlisting}
# Ceci est un titre

Et ceci est un paragraphe.

- Et ceci est une liste à points
- Avec plusieurs points

On peut aussi mettre du texte en gras, comme dans l'exemple : **comme ceci**.
L'italique a _aussi_ une syntaxe.
    \end{lstlisting}
  \end{exampleblock}

  \pause

  Pour afficher des blocks, on utilise la syntax suivante:

  \begin{exampleblock}{Blocs en Slipshow}
    \begin{lstlisting}
{.block}
Ceci est un block


{.block}
> Ceci est un block
>
> Sur plusieurs lignes !
\end{lstlisting}
  \end{exampleblock}

  
  
\end{frame}

\end{document}
